2.2. \underline{Secciones unidad}. Consideremos primero la sección
unidad de $E$ considerada como extensión de $Q_P$ por $G_P$,
\begin{equation}\tag{2.2.1}
 e_{E/P}\in\Gamma(E/P)=\operatorname{Hom}_P(P,E),
\end{equation}
y, simétricamente, la sección unidad de $E$ considerada como extensión
\sisi{de $P_G$ por $G_Q$}{de $P_Q$ por $G_Q$},
\begin{equation}\tag{2.2.2}
 e_{E/Q}\in\Gamma(E/Q)=\operatorname{Hom}_Q(Q,E).
\end{equation}
En particular, \underline{el torsor $E$ sobre $P\times Q$ bajo
$G_{P\times Q}$ queda canónicamente trivializado sobre
\sisi{$1_P\times Q$}{$P\times1_Q$} por $e_{E/P}$, y sobre
\sisi{$P\times1_Q$}{$1_P\times Q$} por $e_{E/Q}$. Estas
trivializaciones coinciden sobre la sección unidad
$1_P\times1_Q$ de $P\times Q$}, y por tanto definen la misma sección

%%%%%%%%%%  scan idx 172  |  folio 161  %%%%%%%%%%

\begin{equation}\tag{2.2.3}
 e_E\in\Gamma(E),
\end{equation}
que, por abuso de lenguaje, se denominará la \underline{sección unidad}
(o \underline{sección neutra}) de $E$.

Para probar la igualdad afirmada, sean $e,e'$ las dos secciones de $E$
deducidas respectivamente de $e_{E/P}$ y $e_{E/Q}$. Se caracterizan
como las secciones unidad de la imagen inversa $E_{1_P,1_Q}$ de $E$ por
la sección unidad de $P\times Q$, respecto de las dos leyes de
composición $\varphi$ y $\psi$, respectivamente. Aplicando (2.1.1.2) a
las cuatro secciones $e,e',e',e$ de $E_{1,1}$, se obtiene $e$ en el
primer miembro y $e'$ en el segundo; por tanto, $e=e'$. Se sigue además
que \underline{las dos leyes de composición $\varphi$ y $\psi$ sobre
$E_{1,1}$ son idénticas}. En efecto, por transporte de estructura
mediante los morfismos $g\mapsto g.e$ y $g\mapsto g.e'$,
respectivamente, ambas corresponden a la ley de grupo de $G$, y
$e=e'$.

La extensión de $Q_P$ por $G_P$ sobre la base $P$, restringida sobre
la sección unidad $1_P$ de $P$, da una extensión de $Q$ por $G$ sobre
la imagen inversa de $1_P\times Q$ en \sisi{$G$}{$E$}. Se comprueba de
inmediato que el homomorfismo $e_{E/Q}$ es multiplicativo para esta
estructura de extensión, por lo que la extensión queda canónicamente
escindida. Simétricamente, la extensión de $P$ por $G$ obtenida al
restringir la extensión de $P_Q$ por $G_Q$ sobre $Q$ a la sección
unidad $1_Q$ queda canónicamente escindida por $e_{E/P}$.

